\section*{第 \ref{ch:g6:matter-from-food} 章——物质来自食物}

\begin{solution}{exo:g6:matter-from-food:1}
\omterm{def:g5:food-webs:roles}{消费者}将所吃的\omterm{def:g6:plants-are-producers:matter}{有机物}转化为自身的
\omterm{def:g6:plants-are-producers:matter}{有机物}。\omterm{def:g5:food-webs:roles}{生产者}则以纯粹的无机原料，借助光能，
制造新的\omterm{def:g6:plants-are-producers:matter}{有机物}。
\end{solution}

\begin{solution}{exo:g6:matter-from-food:2}
\omterm{def:g4:journey-of-food:digestion}{消化}将草分解为\omterm{def:g4:journey-of-food:digestion}{营养物质}
（\cref{ch:g4:journey-of-food}）；\omterm{prop:g4:journey-of-food:absorb}{血液}将其运送
遍布全身（\cref{ch:g4:heart-and-blood}）；身体的\omterm{def:g5:first-look-at-cells:cell}{细胞}
将其重新组装为牛的物质，并随身体生长而分裂
（\cref{ch:g5:first-look-at-cells}）。
\end{solution}

\begin{solution}{exo:g6:matter-from-food:3}
生长（小牛增加的千克数）；更新（脱落的\omterm{def:g1:animals-around-us:covering}{皮毛}、
新生的\omterm{def:g1:animals-around-us:covering}{羽毛}、鹿角）；修复（愈合的\omterm{def:g3:bones-and-muscles:skeleton}{骨骼}）；产物
（奶、蛋、羊毛）。
\end{solution}

\begin{solution}{exo:g6:matter-from-food:4}
未被吸收的——随粪便排出；与吸入的氧气
一同燃烧、用于运动和保暖的——以呼出的气体排出；
其余部分则用于构建生长和产物。
\end{solution}

\begin{solution}{exo:g6:matter-from-food:5}
因为她每天产出的 \qty{25}{kg} 牛奶属于\omterm{def:g6:plants-are-producers:matter}{有机物}，
必须由当天的饲料重新合成；牛奶就是产出，
而产出即是重新合成的食物——没有输入，就没有输出。
\end{solution}

\begin{solution}{exo:g6:matter-from-food:6}
由母鸡所吃谷物中的\omterm{def:g4:journey-of-food:digestion}{营养物质}，在母鸡\emph{体内}组装而成，
但\omterm{def:g6:plants-are-producers:matter}{有机物}本身早在母鸡之前，
就已在谷物植株的叶片中制成。母鸡完成了转化，
\omterm{def:g1:plants-around-us:plant}{植物}完成了生产。
\end{solution}

\begin{solution}{exo:g6:matter-from-food:7}
每个箭头处，所吃物质的大部分被燃烧，用于
维持生命，或随粪便丢失，只有很小一部分
成为捕食者自身的身体——这是下一层唯一
可取食的部分。因此每一层所拥有的物质，
只是其下一层的一小部分：金字塔就这样按算术逐层收窄。
\end{solution}

\begin{solution}{exo:g6:matter-from-food:8}
输入：肉类或猫粮（偶尔还有老鼠）。用于构建的：
幼猫时期的生长、每次换毛时更新的\omterm{def:g1:animals-around-us:covering}{皮毛}、
擦伤后的修复。输出：花园里的粪便，以及燃烧掉的
部分，为每一次扑击、呼噜和温暖的午睡提供\omterm{prop:g3:balanced-diet:jobs}{能量}。
\end{solution}

\begin{solution}{exo:g6:matter-from-food:9}
不同之处在于来源与装备：草与
兔子，研磨器官加庞\omterm{def:g4:journey-of-food:tract}{大肠}道，与捕捉器官加短肠。
相同之处在于交易本身：二者都将所食的\omterm{def:g6:plants-are-producers:matter}{有机物}分解为\omterm{def:g4:journey-of-food:digestion}{营养物质}，
再将其重建为自身的
物质。
\end{solution}

\begin{solution}{exo:g6:matter-from-food:10}
大部分作为燃烧的部分离开身体，用于支持一年中的
奔跑、思考，以及维持体温在
\qty{37}{\degreeCelsius}——另有一部分从未被吸收，
随粪便排出。只有剩下的小部分，
即 \qty{3}{kg}，形成了新的身体：
孩子的账本与奶牛的一样平衡。
\end{solution}

\begin{solution}{exo:g6:matter-from-food:11}
蜗牛的身体从食物和外界\omterm{def:g6:exploring-our-environment:components}{环境}中提取贝壳的
\omterm{def:g6:exploring-our-environment:components}{矿物质}成分，加以运输，并逐层沉积在
自身的结构之中。物质本身是无机的；而\emph{建造}
则是活\omterm{def:g1:animals-around-us:animal}{动物}的工作——
一件由矿物材料制成的造物。
\end{solution}

\begin{solution}{exo:g6:matter-from-food:12}
经由猪转化时，土豆的大部分物质被用于猪自身的
生存——燃烧以取暖和运动，随粪便丢失——
只有很小一部分化为猪肉。若直接食用土豆，
便可省去猪的这些消耗。
这笔差额记在了猪的账上，
而没有端上餐桌。
\end{solution}

\begin{solution}{pb:g6:matter-from-food:1}
\textbf{1.} 牧场的草和菜园的\omterm{def:g1:plants-around-us:plant}{植物}
（以及任何果树）。输入物：土壤中的水和\omterm{prop:g4:what-plants-need:needs}{无机盐}，
空气中的二氧化碳。
\omterm{prop:g3:balanced-diet:jobs}{能量}：光。

\textbf{2.} 在草植株自身的叶片中——每一克都是如此：
\omterm{def:g5:food-webs:roles}{生产者}在原地、白天进行制造。

\textbf{3.} 归还所借用的：粪肥腐熟后
补充菜地土壤中的\omterm{prop:g4:what-plants-need:needs}{无机盐}，
以替代收获时带走的部分。

\textbf{4.} 因为奶牛是\omterm{def:g5:food-webs:roles}{消费者}，
其交易需要\emph{有机}物质来分解和重建。
\omterm{prop:g4:what-plants-need:needs}{无机盐}与水是\omterm{def:g5:food-webs:roles}{生产者}的输入物，
奶牛没有将\omterm{def:g6:exploring-our-environment:components}{矿物质}构建为身体的机制——
这正是她所缺少的交易能力。

\textbf{5.} 一部分未被吸收便排出，即牛粪饼；
很大一部分与吸入的氧气一起燃烧，为行走、
咀嚼和保暖提供\omterm{prop:g3:balanced-diet:jobs}{能量}；其余部分化为牛奶
和牛的身体（账目在当天平衡）。

\textbf{6.} 输入：谷物。建成：卵、她自身的维持，以及——
换羽期间——一整套新\omterm{def:g1:animals-around-us:covering}{羽毛}。输出：粪便与
燃烧掉的部分。换羽期间，“建成”一项用于
\omterm{def:g1:animals-around-us:covering}{羽毛}；输入不变时，卵一项必须
暂时减少。

\textbf{7.} 草 $\rightarrow$ 母鸡（经谷物，谷物本身是
一种\omterm{def:g1:plants-around-us:plant}{植物}产物）$\rightarrow$ 小鸡 $\rightarrow$ 鹰；
或者在牧场一侧：草 $\rightarrow$ 奶牛。物质
只在绿色环节（草、谷类\omterm{def:g1:plants-around-us:plant}{植物}）被\emph{制造}；
每个\omterm{def:g1:animals-around-us:animal}{动物}环节——母鸡、小鸡、鹰、奶牛——只
是对其进行转化。

\textbf{8.} 面包：小麦植株的叶。水煮蛋：母鸡
源自谷物——谷类\omterm{def:g1:plants-around-us:plant}{植物}的叶。牛奶：奶牛
源自草——草的叶。三条路径，三片叶，
都依赖于光。

\textbf{9.} 每一层只传递一部分：在
\qty{10000}{kg} 草中，兔种群仅构建
其中很小的一部分作为自身躯体；在这一部分中，狐狸仅构建
又一部分。两层损耗使牧场的草
的量缩减为一窝狐狸的量——她的估算
正是金字塔式的算术。

\textbf{10.} 若直接食用，这块地的产物
直接到达家庭，中间没有\omterm{def:g5:food-webs:roles}{消费者}层级：其中
没有任何部分在\omterm{def:g1:animals-around-us:animal}{动物}的生命活动中被燃烧，也没有随粪便损失。
而经牧场的途径，要先支付奶牛的全部开销，因此
每平方米能养活的人更少。

\textbf{11.} 草不再生产，因此奶牛的输入
先中断：牛奶先减少（产出是重新构建的食物），
接着她燃烧自身的脂肪储备，体重下降——
“建成”一项转为负值。阳光是\omterm{prop:g3:balanced-diet:jobs}{能量}，而非
物质：没有水，草就没有可用来
构建的原料，因此牧场不能指望天上的阳光来挽救。

\textbf{12.} \omterm{def:g5:food-webs:roles}{生产者}：农场的\omterm{def:g1:plants-around-us:plant}{植物}
以\omterm{prop:g4:what-plants-need:needs}{无机盐}和光为原料，制造其全部\omterm{def:g6:plants-are-producers:matter}{有机物}。\omterm{def:g5:food-webs:roles}{消费者}：
奶牛、母鸡和家人转化这些物质，燃烧其中大部分
并将其余部分构建为自身。缺失的是\omterm{def:g5:food-webs:roles}{分解者}——它们
必须把粪便、牛粪饼和秋季的残体归还为
\omterm{def:g6:exploring-our-environment:components}{矿物质}；它们是下一章的主题。
\end{solution}
